%======================================================================
\section{Deferred proofs}
\label{app:proofs}

\subsection{Proof of \texorpdfstring{\cref{prop:atm}}{the ATM proposition}}
\label{app:atmproof}

The smile is defined implicitly by $\Black\big(k,w(k)\big)=c(k)$.  The
plan: compute the five ATM partials of $\Black$, differentiate the
defining identity twice, solve, and convert to volatility units.
Throughout, $w_0=w(0)$, $d=\sqrt{w_0}/2$, and all partials are evaluated
at $(k,w)=(0,w_0)$.

\paragraph{Step 1: the five Black partials at the money.}
From \eqref{eq:black}, with $d_\pm=-k/\sqrt{w}\pm\sqrt{w}/2$, the
identity $e^{k}\phin(d_-)=\phin(d_+)$ (verified by
$d_+^2-d_-^2=(d_++d_-)(d_+-d_-)=(-2k/\sqrt{w})\cdot\sqrt{w}=-2k$) gives
the standard cancellations:
\begin{align*}
\partial_k\Black
&=\phin(d_+)\cdot\big(-\tfrac{1}{\sqrt w}\big)
 -e^{k}\PhiN(d_-)-e^{k}\phin(d_-)\cdot\big(-\tfrac{1}{\sqrt w}\big)
 =-e^{k}\PhiN(d_-),\\
\partial_w\Black
&=\phin(d_+)\,\partial_w d_+-e^{k}\phin(d_-)\,\partial_w d_-
 =\phin(d_+)\,\partial_w(d_+-d_-)
 =\frac{\phin(d_+)}{2\sqrt{w}} .
\end{align*}
Differentiating these once more, using $\phin'(t)=-t\,\phin(t)$:
\begin{align*}
\partial_{kk}\Black
&=-e^{k}\PhiN(d_-)-e^{k}\phin(d_-)\,\partial_k d_-
 =-e^{k}\PhiN(d_-)+\frac{\phin(d_+)}{\sqrt{w}},\\
\partial_{kw}\Black
&=\partial_k\!\left[\frac{\phin(d_+)}{2\sqrt w}\right]
 =\frac{-d_+\phin(d_+)}{2\sqrt w}\cdot\big(-\tfrac{1}{\sqrt w}\big)
 =\frac{d_+\,\phin(d_+)}{2w},\\
\partial_{ww}\Black
&=\frac{-d_+\phin(d_+)\,\partial_w d_+}{2\sqrt w}
 -\frac{\phin(d_+)}{4w^{3/2}},
\qquad
\partial_w d_+=\frac{k}{2w^{3/2}}+\frac{1}{4\sqrt w}.
\end{align*}
At $k=0$: $d_+=d$, $d_-=-d$, $\PhiN(d_-)=1-\PhiN(d)$,
$\partial_wd_+=1/(4\sqrt{w_0})$, so
\begin{equation}
\begin{aligned}
\partial_k\Black&=-\big(1-\PhiN(d)\big), &
\partial_w\Black&=\frac{\phin(d)}{2\sqrt{w_0}}, &
\partial_{kk}\Black&=-\big(1-\PhiN(d)\big)+\frac{\phin(d)}{\sqrt{w_0}},\\
\partial_{kw}\Black&=\frac{\phin(d)}{4\sqrt{w_0}}, &
\partial_{ww}\Black&=-\frac{\phin(d)}{4w_0^{3/2}}
-\frac{\phin(d)}{16\sqrt{w_0}} . &&
\end{aligned}
\label{eq:blackpartials}
\end{equation}

\paragraph{Step 2: first derivative.}
Differentiating $\Black(k,w(k))=c(k)$ once and solving,
\[
w'(0)=\frac{c'(0)-\partial_k\Black}{\partial_w\Black}
=\frac{-(1-u_*)+(1-\PhiN(d))}{\phin(d)/(2\sqrt{w_0})}
=\frac{2\sqrt{w_0}}{\phin(d)}\,\big(u_*-\PhiN(d)\big)
=\frac{2\sqrt{w_0}}{\phin(d)}\,\Delta,
\]
which is \eqref{eq:atm_w}.

\paragraph{Step 3: second derivative.}
Differentiating twice,
\[
\partial_{kk}\Black+2\,\partial_{kw}\Black\,w'(0)
+\partial_{ww}\Black\,w'(0)^{2}+\partial_w\Black\,w''(0)=c''(0),
\]
and solving for $w''(0)$ with $c''(0)=f_*-(1-u_*)$ and
\eqref{eq:blackpartials}:
\begin{align*}
w''(0)
&=\frac{2\sqrt{w_0}}{\phin(d)}
\left[f_*-(1-u_*)+\big(1-\PhiN(d)\big)-\frac{\phin(d)}{\sqrt{w_0}}
-\frac{\phin(d)}{2\sqrt{w_0}}\,w'(0)
+\Big(\frac{\phin(d)}{4w_0^{3/2}}+\frac{\phin(d)}{16\sqrt{w_0}}\Big)w'(0)^{2}
\right]\\
&=\frac{2\sqrt{w_0}}{\phin(d)}\,f_*
+\underbrace{\frac{2\sqrt{w_0}}{\phin(d)}\,\Delta}_{=\,w'(0)}
-\;2\;-\;w'(0)
+\Big(\frac{1}{2w_0}+\frac18\Big)w'(0)^{2}
=\frac{2\sqrt{w_0}}{\phin(d)}\,f_*-2
+\Big(\frac{1}{2w_0}+\frac18\Big)w'(0)^{2},
\end{align*}
the promised \eqref{eq:atm_wpp}: the digital-mismatch terms cancel
against the cross term exactly.

\paragraph{Step 4: volatility units.}
With $\sigma(k)=\sqrt{w(k)/\tau}$: $\sigma'=w'/(2\sqrt{\tau w})$ and
$\sigma''=w''/(2\sqrt{\tau w})-w'^{2}/(4\sqrt{\tau}\,w^{3/2})$.  At
$k=0$, substituting $w'(0)=2\sqrt{w_0}\Delta/\phin(d)$:
\[
\sigma'(0)=\frac{\Delta}{\phin(d)\sqrt{\tau}},
\qquad
\sigma''(0)
=\frac{1}{\sqrt\tau}\left[
\frac{f_*}{\phin(d)}-\frac{1}{\sqrt{w_0}}
+\Big(\frac{1}{2w_0}+\frac18\Big)\frac{w'(0)^2}{2\sqrt{w_0}}
-\frac{w'(0)^{2}}{4w_0^{3/2}}
\right],
\]
and the bracket's last two terms combine:
$\big(\tfrac{1}{2w_0}+\tfrac18\big)\tfrac{1}{2\sqrt{w_0}}
-\tfrac{1}{4w_0^{3/2}}
=\tfrac{1}{4w_0^{3/2}}+\tfrac{1}{16\sqrt{w_0}}-\tfrac{1}{4w_0^{3/2}}
=\tfrac{1}{16\sqrt{w_0}}$, giving with
$w'(0)^{2}=4w_0\Delta^{2}/\phin(d)^{2}$ the final form
\[
\sigma''(0)=\frac{1}{\sqrt{\tau}}
\left[\frac{f_*}{\phin(d)}-\frac{1}{\sqrt{w_0}}
+\frac{\sqrt{w_0}}{4}\Big(\frac{\Delta}{\phin(d)}\Big)^{2}\right],
\]
which is \eqref{eq:atm_sigma}. \hfill$\square$

\subsection{Proof of \texorpdfstring{\cref{prop:universal}}{the universality proposition}}
\label{app:universalproof}

\paragraph{Setup.}
Let $X^{\ast}$ satisfy the hypothesis: quantile
$Q^{\ast}(u)=x^{\ast}(\logit u)$ with
$x^{\ast\prime}(z)=e^{g^{\ast}(\Lambda(z))}$ for a continuous
$g^{\ast}\colon[0,1]\to\R$ with $g^{\ast}(1)<0$, and $\E[e^{X^{\ast}}]=1$.
Write $\lambda_\pm^{\ast}$ for its endpoint scales and fix
$\bar\lambda\in(\lambda_+^{\ast},1)$.

\paragraph{Step 1: polynomial approximation inside the family.}
By the Weierstrass approximation theorem, choose polynomials $p_N$ of
degree at most $N$ with
$\varepsilon_N:=\|p_N-g^{\ast}\|_{\infty}\to0$.  Every polynomial of
degree $\le N$ is exactly representable in the basis of \eqref{eq:lqd}
(the endpoint chord spans degrees $0,1$ and $P_2,\dots,P_N$ the rest),
so $p_N$ defines a candidate slice with coefficients $\theta_N$.
Admissibility: its right endpoint value is
$p_N(1)\le g^{\ast}(1)+\varepsilon_N<\log\bar\lambda<0$ for all $N$ past
some $N_0$ --- admissibility is an open condition, so uniform
approximation lands inside it.  Fix such $N$ hereafter; write $g_N=p_N$,
and let $\bar{x}_N$, $m_N$, $x_N$, $Q_N$, $c_N$ be the associated
objects.

\paragraph{Step 2: transports converge, with a uniform tail envelope.}
Pointwise, $|e^{g_N}-e^{g^{\ast}}|\le
e^{g^{\ast}}\big(e^{\varepsilon_N}-1\big)$, so for every $z$,
\[
\big|\bar{x}_N(z)-\bar{x}^{\ast}(z)\big|
\le\big(e^{\varepsilon_N}-1\big)\int_{0}^{|z|}
e^{g^{\ast}(\Lambda(\pm t))}\dd t
\le\big(e^{\varepsilon_N}-1\big)\,e^{\|g^{\ast}\|_\infty}\,|z| .
\]
For the envelope, no derivative information about $g_N$ is needed ---
sup-norm closeness suffices.  Since $g^{\ast}$ is continuous with
$g^{\ast}(1)=\log\lambda_+^{\ast}<\log\bar\lambda$, pick $z_0$ and
$\delta>0$ with $g^{\ast}(\Lambda(t))\le\log\bar\lambda-\delta$ for all
$t\ge z_0$; then for all $N$ with $\varepsilon_N<\delta$,
$e^{g_N(\Lambda(t))}\le\bar\lambda$ on $[z_0,\infty)$, so integrating
from $z_0$,
\begin{equation}
e^{\bar{x}_N(z)}\;\le\;C_1\,e^{\bar\lambda z}\quad(z\ge0),
\qquad
e^{\bar{x}_N(z)}\;\le\;C_1\quad(z\le0),
\label{eq:envelopeB}
\end{equation}
with $C_1=\exp\big(z_0\,e^{\|g^{\ast}\|_\infty+\delta}\big)$
independent of $N$ (the $z\le0$ bound holds because $\bar{x}_N$ is
increasing with $\bar{x}_N(0)=0$).

\paragraph{Step 3: normalizers and shifts converge.}
$I_N=\int e^{\bar{x}_N}\rho\,\dd z$.  By Step 2 the integrands converge
pointwise to $e^{\bar{x}^{\ast}}\rho$ and are dominated by
$C_1e^{\bar\lambda z}\rho(z)\le C_1e^{(\bar\lambda-1)z}$ on the right and
$C_1e^{z}$ on the left, both integrable; dominated convergence gives
$I_N\to I^{\ast}$ and $m_N=-\log I_N\to m^{\ast}$.  Hence
$x_N\to x^{\ast}$ uniformly on every compact $z$-interval, i.e.\
$Q_N\to Q^{\ast}$ uniformly on compact subsets of $(0,1)$.

\paragraph{Step 4: calls converge uniformly.}
Couple all the laws by one uniform rank: $Y_N=e^{Q_N(U)}$,
$Y^{\ast}=e^{Q^{\ast}(U)}$.  For every $k$,
\[
\big|c_N(k)-c^{\ast}(k)\big|
=\Big|\E\big[(Y_N-e^{k})^{+}-(Y^{\ast}-e^{k})^{+}\big]\Big|
\le\E\big|Y_N-Y^{\ast}\big|,
\]
since $t\mapsto(t-e^k)^+$ is $1$-Lipschitz --- one bound, uniform in
$k$.  Now $\E|Y_N-Y^{\ast}|=\int_0^1|e^{Q_N}-e^{Q^{\ast}}|\dd u\to0$: the
integrand tends to zero pointwise (Step 3) and is dominated by
$e^{Q_N}+e^{Q^{\ast}}\le 2C_2(1-u)^{-\bar\lambda}+2C_2$, integrable
because $\bar\lambda<1$.  Here $C_2$ --- the only place uniform control
of the shifts is needed --- is
$C_2=C_1\,e^{\sup_N|m_N|}\cdot2^{\bar\lambda}$: the factor $C_1$ from
\eqref{eq:envelopeB}, the shift factor finite because $m_N\to
m^{\ast}$ (Step 3) so $\sup_N|m_N|<\infty$, and the last factor from
translating the envelope through the two-sided bound
$\tfrac12e^{-z}\le1-\Lambda(z)\le e^{-z}$ for $z\ge0$.

\paragraph{Step 5: Wasserstein convergence of the return laws.}
For laws on $\R$ with finite $p$-th moments, $p\ge1$, the quantile
coupling is optimal:
$W_p(Y_N,Y^{\ast})^{p}=\int_0^1|e^{Q_N}-e^{Q^{\ast}}|^{p}\dd u$
\citep[Ch.~2]{Villani2003}.  For $1\le p<1/\lambda_+^{\ast}$ (a nonempty
range, since $\lambda_+^{\ast}<1$) choose
$\bar\lambda\in(\lambda_+^{\ast},\min(1,1/p))$ in Step 2; then the
integrand is dominated by $2^{p}C_2^{p}\big((1-u)^{-p\bar\lambda}+1\big)$
with $p\bar\lambda<1$, integrable, and pointwise convergence again
finishes it.  (For $p>1/\lambda_+^{\ast}$ the target law itself lacks
the moment, so no metric statement is attempted.)

\paragraph{Step 6: exclusions.}
At any finite order, $x_N'$ is continuous and bounded between two
positive constants on every compact $z$-set; hence $Q_N$ is a strictly
increasing homeomorphism of $(0,1)$ onto an unbounded interval.  An atom
of the law is a jump of the quantile function; a zero-density gap is a
flat of it; bounded support means $Q$ bounded; default mass at zero
means $Q=-\infty$ on a rank interval --- each impossible for such
$Q_N$.  A Gaussian (or any super-exponential) upper tail forces
$g(u)\to-\infty$ as $u\to1$ (\cref{sec:heuristic}), impossible for a
polynomial.

For the closure dichotomy, let $g_N\to g_\infty$ uniformly on $[0,1]$
with every $g_N$ admissible; $g_\infty$ is continuous and
$g_\infty(1)=\lim_N g_N(1)\le0$.

\emph{Branch $g_\infty(1)<0$.}  Steps 2--4 apply verbatim with
$g_\infty$ as the target: normalizers, shifts, and quantiles converge
(locally uniformly), the call curves converge uniformly --- pointwise
quantile convergence gives weak convergence --- and the limit law is of
the hypothesis class: positive continuous density, exponential tails,
admissible scales.

\emph{Branch $g_\infty(1)=0$.}  Weak convergence can fail, and the
paper's own solvable slice shows it: take $g_N\equiv\log(1-1/N)$,
constant speeds approaching the wall.  By \eqref{eq:mgf},
$I_N=\pi s_N/\sin(\pi s_N)$ with $s_N=1-1/N$, so
$I_N=\pi(1-1/N)/\sin(\pi/N)\to\infty$, hence $m_N=-\log I_N\to-\infty$
and every quantile $Q_N(u)=m_N+s_N\logit u\to-\infty$: for any fixed
level $q$, $\Prob(X_N\le q)\to1$, the mass escapes to $-\infty$, and no
weak limit exists.  Suppose instead that the laws \emph{do} converge
weakly to a probability law $\mu$.  First, $\{m_N\}$ is bounded: above,
because $I_N\ge\int_{-1}^{1}e^{\bar x_N}\rho\,\dd z$ is bounded away
from zero by locally uniform convergence of $\bar x_N$; and below,
because $m_{N_j}\to-\infty$ along a subsequence would give
$Q_{N_j}(u)=m_{N_j}+\bar x_{N_j}(\logit u)\to-\infty$ at every rank,
i.e.\ $\Prob(X_{N_j}\le q)\to1$ for every $q$, contradicting the
tightness of a weakly convergent sequence.  Next, any subsequential
limit $\bar m$ of $m_N$ gives $Q_N\to\bar m+\bar x_\infty\circ\logit$
pointwise on $(0,1)$ along that subsequence, and pointwise quantile
convergence (at continuity points --- here everywhere) is equivalent to
weak convergence on $\R$; so $\mu$'s quantile equals
$\bar m+\bar x_\infty\circ\logit$, which pins $\bar m$ and shows the
whole sequence $m_N$ converges.  The limit quantile is $C^{1}$ in
log-odds with derivative $e^{g_\infty(\Lambda(z))}$, bounded below by
$e^{\min_{[0,1]}g_\infty}>0$; hence $\mu$ has support all of $\R$ and,
by the computation \eqref{eq:density} with $g_\infty$ in place of $g$,
a strictly positive continuous density $u(1-u)e^{-g_\infty(u)}$ at rank
$u$ --- no atom, no gap, no bounded support, and no Gaussian-or-faster
tail (which would force $g_\infty(u)\to-\infty$ at an endpoint).  Its
right tail sits at the wall: for every $\delta>0$ there is $z_\delta$
with $e^{g_\infty(\Lambda(z))}\ge1-\delta$ for $z\ge z_\delta$, so the
quantile grows at asymptotic rate at least $1-\delta$, and the
moment-strip computation of \cref{prop:strip} (with $1-\delta$ in place
of $\lambda_+$, $\delta$ arbitrary) gives $\E_\mu[e^{rX}]=\infty$ for
every $r>1$; by Fatou, $\E_\mu[e^{X}]\le\liminf_N\E[e^{X_N}]=1$, with
strict inequality possible.  The boundary limit, when it exists, is
thus not a member of the class of (i) --- its tail scale has been lost
--- but neither is it any excluded law.  In both branches, no excluded
law is reached. \hfill$\square$

\begin{remark}
Step 4's one-line Lipschitz bound is the reason the paper prices
through the rank coupling everywhere: the same trick gives the
stability of every price-linear functional under parameter
perturbations, and it is the continuous heart of the delta-method
bands mentioned in \cref{app:implementation}.
\end{remark}
